\SourcePage{713}{716}
\section{Elastic scattering in the presence of inelastic processes}

Collisions accompanied by a change in the internal state of the colliding
particles are called inelastic. Here these changes are understood in the
broadest sense; even the type of particle may change. Thus atoms may be
excited or ionized, and nuclei may be excited or disintegrated. When a
collision, for example a nuclear reaction, may be accompanied by different
physical processes, one speaks of different \emph{reaction channels}.

The presence of inelastic channels also has a definite effect on the
properties of elastic scattering.

In the general case of several reaction channels, the asymptotic wave
function of the colliding-particle system is a sum with one term for each
possible channel. Among these is a term describing the particles in their
unchanged initial state, the \emph{entrance channel}. It is the product of the
wave functions of the particles' internal states and a function describing
their relative motion in the center-of-mass system. This last function is of
interest here; denote it by $\psi$ and determine its asymptotic form.

In the entrance channel, $\psi$ consists of an incident plane wave and an
outgoing spherical wave corresponding to elastic scattering. It may also be
represented as a sum of incoming and outgoing waves, as in \S~123. The
difference is that the asymptotic expression for the radial functions
$R_l(r)$ cannot be taken as a standing wave. A standing wave is the sum of
incoming and outgoing waves of equal amplitude. This accords with the
physical meaning of purely elastic scattering, but when inelastic channels
are present, the outgoing-wave amplitude must be smaller than the
incoming-wave amplitude. Thus the asymptotic expression for $\psi$ is given
by (123.9),

\SourcePage{714}{717}
\begin{equation}
 \psi=\frac1{2ikr}\sum_{l=0}^{\infty}(2l+1)P_l(\cos\theta)
 \left[(-1)^{l+1}e^{-ikr}+S_le^{ikr}\right],
 \tag{142.1}
\end{equation}
except that $S_l$ is no longer defined by (123.10), but is a generally
complex quantity with modulus less than unity. The elastic scattering
amplitude is expressed through these quantities by (123.11):
\begin{equation}
 f(\theta)=\frac1{2ik}\sum_{l=0}^{\infty}(2l+1)(S_l-1)P_l(\cos\theta).
 \tag{142.2}
\end{equation}
In place of (123.12), the total elastic scattering cross-section $\sigma_e$ is
\begin{equation}
 \sigma_e=\frac{\pi}{k^2}\sum_{l=0}^{\infty}(2l+1)|1-S_l|^2.
 \tag{142.3}
\end{equation}

The total inelastic scattering cross-section, also called the \emph{reaction
cross-section} $\sigma_r$, summed over every possible channel, can likewise be
expressed in terms of $S_l$. For each $l$, the outgoing-wave intensity is
reduced relative to the incoming-wave intensity by the factor $|S_l|^2$.
This reduction must be attributed entirely to inelastic scattering. Hence
\begin{equation}
 \sigma_r=\frac{\pi}{k^2}\sum_{l=0}^{\infty}(2l+1)(1-|S_l|^2),
 \tag{142.4}
\end{equation}
and the total cross-section is
\begin{equation}
 \sigma_t=\sigma_e+\sigma_r
 =\frac{2\pi}{k^2}\sum_{l=0}^{\infty}(2l+1)(1-\operatorname{Re}S_l).
 \tag{142.5}
\end{equation}

The partial elastic scattering amplitude with angular momentum $l$, defined
by (123.15), is
\begin{equation}
 f_l=\frac{S_l-1}{2ik},
 \tag{142.6}
\end{equation}

\SourcePage{715}{718}
and each term of the sums in (142.3), (142.4) is the partial elastic or
inelastic cross-section for particles of angular momentum $l$:
\begin{equation}
 \begin{aligned}
 \sigma_e^{(l)}&=\frac{\pi}{k^2}(2l+1)|1-S_l|^2,\\
 \sigma_r^{(l)}&=\frac{\pi}{k^2}(2l+1)(1-|S_l|^2),\\
 \sigma_t^{(l)}&=\frac{2\pi}{k^2}(2l+1)(1-\operatorname{Re}S_l).
 \end{aligned}
 \tag{142.7}
\end{equation}

The value $S_l=1$ means that there is no scattering at all for this $l$. The
case $S_l=0$ corresponds to complete ``absorption'' of particles with angular
momentum $l$: the partial outgoing wave with this $l$ is absent from (142.1).
The elastic and inelastic cross-sections are then equal:
\begin{equation}
 \sigma_e^{(l)}=\sigma_r^{(l)}=\frac{\pi}{k^2}(2l+1).
 \tag{142.8}
\end{equation}
We also note that elastic scattering may occur without inelastic scattering,
when $|S_l|=1$, but the converse is impossible: inelastic scattering
necessarily entails elastic scattering. For a specified
$\sigma_r^{(l)}$, the partial elastic cross-section must lie in the interval
\begin{equation}
 \sqrt{\sigma_0}-\sqrt{\sigma_0-\sigma_r^{(l)}}\leq
 \sqrt{\sigma_e^{(l)}}\leq
 \sqrt{\sigma_0}+\sqrt{\sigma_0-\sigma_r^{(l)}},
 \tag{142.9}
\end{equation}
where $\sigma_0=(2l+1)\pi/k^2$.

Taking $f(\theta)$ from (142.2) at $\theta=0$ and comparing it with (142.5)
gives
\begin{equation}
 \operatorname{Im}f(0)=\frac{k}{4\pi}\sigma_t,
 \tag{142.10}
\end{equation}
which generalizes the optical theorem (125.9). Here $f(0)$ is still the
elastic zero-angle scattering amplitude, but $\sigma_t$ includes the
inelastic part as well.

The imaginary parts of the partial amplitudes $f_l$ are related to
$\sigma_t^{(l)}$ by
\begin{equation}
 \operatorname{Im}f_l=\frac{k}{4\pi}\frac{\sigma_t^{(l)}}{2l+1},
 \tag{142.11}
\end{equation}
which follows directly from (142.6), (142.7).

\SourcePage{716}{719}
The fact that the coefficients $S_l$ in the asymptotic wave function do not
have unit modulus has no effect on the conclusions of \S~128 concerning the
singular points of the elastic scattering amplitude as a function of complex
$E$; they remain valid in the presence of inelastic processes. The analytic
properties do change in that the amplitude is now nonreal on the negative
real semiaxis ($E<0$), and its values on the upper and lower edges of the cut
for $E>0$ are not complex conjugates. More generally, its values at points in
the upper and lower half-planes symmetric about the real axis are not complex
conjugates either.

On passing from the upper edge of the cut to the lower by a full circuit
around $E=0$, the root $\sqrt E$ changes sign; hence the real quantity $k$ at
$E>0$ changes sign. The incoming and outgoing waves in (142.1) exchange
roles, so the new coefficient $S_l$ is $1/S_l$, the reciprocal of its former
value, which is by no means $S_l^*$. Denote the amplitudes on the upper and
lower edges by $f_l(k)$ and $f_l(-k)$; only $f_l(k)$ is, of course, the
physical amplitude. From (142.6),
\[
 f_l(k)=\frac{S_l-1}{2ik},\qquad
 f_l(-k)=-\frac{1/S_l-1}{2ik}.
\]
Eliminating $S_l$ gives
\begin{equation}
 f_l(k)-f_l(-k)=2ikf_l(k)f_l(-k).
 \tag{142.12}
\end{equation}
In the absence of inelastic processes, one would have
$f(-k)=f^*(k)$, and (142.12) and (142.11) would coincide.

Rewriting (142.12) as
\[
 \frac1{f_l(k)}-\frac1{f_l(-k)}=-2ik,
\]
shows that $1/f_l(k)+ik$ must be an even function of $k$. Denoting this
function by $g_l(k^2)$, we have
\begin{equation}
 f_l(k)=\frac1{g_l(k^2)-ik}.
 \tag{142.13}
\end{equation}

\SourcePage{717}{720}
The even function $g_l(k^2)$ is no longer real, however, as it was in
(125.15)\footnote{This argument, and hence the conclusion that $g_l$ is even,
assumes that the interaction decreases sufficiently rapidly as $r\to\infty$
to ensure the absence of cuts in the left half-plane of $E$, thereby allowing
a full circuit around $E=0$.}.

When a particle beam passes through a scattering medium containing many
scattering centers, it gradually weakens as particles undergoing various
collision processes leave the beam. This attenuation is determined entirely
by the elastic zero-angle scattering amplitude and, under certain conditions
(see below), can be described by the following formal method\footnote{The
ideas presented below apply, in particular, to the scattering by nuclei of
fast neutrons with energies of order hundreds of MeV. Their wavelength is so
short that the nucleus may be regarded as an inhomogeneous macroscopic medium.}.

Let $f(0,E)$ be the zero-angle scattering amplitude for each particle of the
medium. Assume that $f$ is small compared with the mean distance
$d\sim(V/N)^{1/3}$ between particles, so each may be treated as a separate
scatterer. Introduce as an auxiliary quantity an effective field
$U_{\mathrm{eff}}$ of a fixed center, defined so that its Born zero-angle
scattering amplitude is precisely the true amplitude $f(0,E)$. This does not
imply that the Born approximation applies to the calculation of $f(0,E)$ from
the true particle interaction. By definition (see (126.4)),
\begin{equation}
 \int U_{\mathrm{eff}}\,dV=-\frac{2\pi\hbar^2}{m}f(0,E),
 \tag{142.14}
\end{equation}
where $m$ is the mass of the scattered particle. Since $f$ is complex, so is
the field thus defined. An estimate of both sides of (142.14) relates its
range $a$ to its magnitude:
\begin{equation}
 a^3U_{\mathrm{eff}}\sim\hbar^2f/m.
 \tag{142.15}
\end{equation}
The definition (142.14) is, of course, not unique. Impose the additional
condition that $U_{\mathrm{eff}}$ satisfy the applicability condition for
perturbation theory:
\begin{equation}
 |U_{\mathrm{eff}}|\ll\hbar^2/ma^2
 \tag{142.16}
\end{equation}

\SourcePage{718}{721}
(then $|f|\ll a$). The attenuation of the scattered beam can then be
described as propagation of a plane wave through a homogeneous medium in
which the particle has the constant potential energy
\begin{equation}
 \overline{U}_{\mathrm{eff}}=\frac NV\int U_{\mathrm{eff}}\,dV
 =-\frac NV\frac{2\pi\hbar^2}{m}f(0,E),
 \tag{142.17}
\end{equation}
obtained by averaging the effective fields of all $N$ particles over the
volume $V$ of the medium. This is clear if one first considers scattering in
a region of the medium which contains many scattering centers but in which
the scattering effect is still small; such regions can be selected by virtue
of (142.16). The attenuation on passing through this region is determined by
the zero-angle amplitude, which in the Born approximation is in turn
determined by the integral of the scattering field over the whole volume of
the region. This means precisely that the scattering properties of interest
are determined completely by the field (142.17) averaged over the volume.

The beam passing through the medium may therefore be described by a plane
wave proportional to $e^{ikz}$ with wave vector
\[
 k=\frac1\hbar\sqrt{2m(E-\overline{U}_{\mathrm{eff}})}.
\]
Introduce the incident-particle wave vector $k_0=\sqrt{2mE}/\hbar$ and write
$k=nk_0$. The quantity
\begin{equation}
 n=\sqrt{1-\frac{\overline{U}_{\mathrm{eff}}}{E}}
 =\sqrt{1+\frac NV\frac{2\pi\hbar^2}{mE}f(0,E)}
 \tag{142.18}
\end{equation}
acts as the \emph{refractive index} of the medium for the beam passing through
it. It is generally complex, since the amplitude is complex, and its imaginary
part determines the attenuation of the beam intensity. If
$E\gg|\overline{U}_{\mathrm{eff}}|$, then (142.18) gives, as it should,
\[
 \operatorname{Im}n=\frac NV\frac{\pi\hbar^2}{mE}\operatorname{Im}f(0,E)
 =\frac NV\frac{\sigma_t}{2k},
\]
where $\sigma_t$ is the total scattering cross-section; the optical theorem
(142.10) has been used. This expression gives the evident attenuation law
\[
 |e^{ikz}|^2\sim\exp\left(-\frac NV\sigma_tz\right).
\]

\SourcePage{719}{722}
Besides absorption, the complex refractive index (142.18), through its real
part, determines the refraction law at entry into and exit from the scattering
medium\footnote{An interesting application of (142.17) is the shift of the
higher levels of an alkali-metal atom immersed in a foreign gas. In a highly
excited state the valence electron is at a mean distance $r$ from the atomic
center large compared with the dimensions $a$ of both the ionic core and the
foreign neutral atoms. The latter atoms within a sphere of radius of order
$r$ act as scattering centers for the valence electron and shift its energy
level by (142.17). Since the de Broglie wavelength of the excited valence
electron is also large compared with $a$, the amplitude is
$f(0,E)\approx-a$, where $a$ is the scattering length (compare (132.9)). The
result is a constant level shift $2\pi\hbar^2a\nu/m$, where $m$ is the
electron mass and $\nu$ is the number density of the foreign gas
(E.~Fermi, 1934).}.

\ProblemHeading

Neutrons are scattered by a heavy nucleus, and their wavelength is small
compared with the nuclear radius $a$ ($ka\gg1$). Assume that every neutron
incident with orbital angular momentum $l<ka\equiv l_0$, that is, with impact
parameter $\rho=\hbar l/mv=l/k<a$, is absorbed by the nucleus, while for
$l>l_0$ it does not interact with the nucleus at all. Determine the elastic
scattering cross-section at small angles.

\SolutionHeading
Under these conditions the neutron motion is mainly semiclassical, while the
elastic scattering results from a small deflection entirely analogous to
Fraunhofer diffraction of light by a black sphere. The required cross-section
can therefore be written directly from the familiar solution of the
diffraction problem\footnote{See II, \S~61, Problem 3. Diffraction by a black
sphere is equivalent to diffraction by a circular aperture in an opaque
screen. The scattering cross-section is obtained by dividing the intensity
of the diffracted waves by the incident current density.}:
\[
 d\sigma_e=\pi a^2\frac{J_1^2(ka\theta)}{\pi\theta^2}\,d\Omega.
\]
The same result follows from (142.3). By assumption, $S_l=0$ for $l<l_0$ and
$S_l=1$ for $l>l_0$. Hence the elastic scattering amplitude is
\[
 f(\theta)=-\frac1{2ik}\sum_{l=0}^{l_0}(2l+1)P_l(\cos\theta).
\]
The main contribution comes from large $l$. Accordingly, replace $2l+1$ by
$2l$, use the approximation (49.6) for $P_l(\cos\theta)$ at small $\theta$,
and replace summation by integration:
\[
 f(\theta)=\frac{i}{k}\int_0^{l_0}lJ_0(\theta l)\,dl
 =\frac{i}{k\theta}l_0J_1(\theta l_0)
 =\frac{ia}{\theta}J_1(ka\theta),
\]

\SourcePage{720}{723}
as required\footnote{The diffraction scattering of fast charged particles by
a ``black'' nucleus may be treated similarly. The boundary value $l_0$ must
then be found by requiring the shortest distance between the nucleus and a
particle moving on a classical trajectory in the Coulomb field to equal the
nuclear radius. One must again put $S_l=0$ for $l<l_0$, while for $l>l_0$,
$S_l=e^{2i\delta_l}$, where $\delta_l$ are the Coulomb phases from (135.11).
See A.~I.~Akhiezer and I.~Ya.~Pomeranchuk, \emph{Some Problems of Nuclear
Theory}, Moscow: Gostekhizdat, 1950, \S~22; J. Physics USSR, 1945, Vol.~9,
p.~471.}. The total elastic scattering cross-section is
\[
 \sigma_e=\pi a^2\int_0^\infty\frac{J_1^2(ka\theta)}{\pi\theta^2}
 \,2\pi\theta\,d\theta=\pi a^2
\]
(the upper integration limit may be extended to infinity because the integral
converges rapidly). As required under these conditions (compare (142.8)), it
equals the absorption cross-section, which is simply the geometric
cross-sectional area of the sphere. The total cross-section is
$\sigma_t=2\pi a^2$.
